We review existing approaches for constructing video query benchmarks, categorizing them into three primary workflows: manual annotation, (M)LLM-based generation, and simulation-based synthesis.
% 我们回顾了构建视频查询基准的现有方法，并将其分为三种主要工作流程：手工标注、基于大语言模型的生成以及基于模拟的合成。
\subsection{Manual Annotation Pipelines}
Early benchmarks primarily rely on crowdsourcing or expert annotation to ensure high semantic quality. A representative work, MedVidQA \cite{MedVid}, proposes a pipeline that retrieves medical instructional videos based on textual queries from trusted health resources (e.g., WikiHow) and employs human annotators to manually label the start and end timestamps of visual segments answering specific questions.
% 早期基准测试主要依赖众包或专家标注来确保高语义质量。一项具有代表性的工作 MedVidQA [1] 提出了一种流程，根据受信任的健康资源（如 MedlinePlus）中的文本查询检索医疗教学视频，并聘请人工标注者手动标记回答特定问题的视觉片段的起止时间戳。
While manual annotation ensures the authenticity of the video source, it suffers from prohibitive scalability costs. As noted in a similar manual annotation work, M$^3$-Med \cite{m3med}, ensuring high-quality temporal localization requires rigorous expert verification, making it impossible to scale to the millions of training examples needed for modern deep learning models. Furthermore, human annotators often struggle to consistently define precise boundaries for fine-grained actions, leading to subjective bias in the ground truth.
% 虽然人工标注可以确保视频来源的真实性，但在可扩展性方面成本高昂。如 M3Med 所指出，确保高质量的时间定位需要严格的专家验证，这使得无法扩展到现代深度学习模型所需的数百万训练样本。此外，人类标注者通常难以为细粒度动作一致地定义准确的边界，从而导致标注结果中存在主观偏差。

\subsection{(M)LLM-based Automated Generation}
To mitigate the scalability constraints inherent in manual annotation, recent research has pivoted towards utilizing Large Language Models (LLMs) or Multi-modal Large Language Models (MLLMs) for automated dataset synthesis.
% 为了缓解人工标注固有的可扩展性限制，近期的研究转向利用大型语言模型（LLM）进行自动化数据集合成。
For instance, HealthVidQA \cite{HealthVid} proposes a framework that leverages domain-adapted LLMs to derive question-answer pairs directly from Automatic Speech Recognition (ASR) transcripts, enabling the construction of large-scale datasets with minimal human intervention.
% 例如，HealthVidQA \cite{HealthVid} 提出了一种框架，该框架利用领域自适应的 LLM 直接从自动语音识别（ASR）转录文本中提取问答对，从而在最大限度减少人工干预的情况下构建大规模数据集。
However, the primary limitation of such text-centric approaches is the inherent "modality gap": by synthesizing queries based solely on textual metadata rather than raw visual signals, these methods frequently fail to ensure precise visual grounding.
% 然而，此类以文本为中心的方法的主要局限在于其固有的“模态鸿沟”：由于这些方法仅基于文本元数据而非原始视觉信号合成查询，因此往往无法确保精确的视觉对应关系。
To address this misalignment, recent efforts like ShareGPT4Video \cite{ShareGPT4Video} employ MLLMs, such as GPT-4V, to generate queries directly from sampled video frames, aiming to better align textual inquiries with visual content.
% 为了解决这种错位问题，ShareGPT4Video \cite{ShareGPT4Video} 等近期项目采用了多模态 LLM（MLLM），例如 GPT-4V，直接从采样的视频帧生成查询，旨在更好地将文本查询与视觉内容对齐。
Nevertheless, these end-to-end generative models operate as probabilistic "black boxes," making them inherently prone to hallucinations—a critical vulnerability in high-stakes medical domains where data integrity is paramount.
% 然而，这些端到端的生成模型如同概率性的“黑箱”，本质上容易产生幻觉——这在数据完整性至关重要的高风险医疗领域是一个关键的漏洞。
Specifically, such models may hallucinate entities not present in the visual stream or synthesize procedural sequences that violate strict medical logic. Correcting these errors necessitates complex, unreliable post-hoc filtering pipelines, thereby undermining the efficiency and trustworthiness of the automated generation process.
% 具体而言，此类模型可能会产生视觉流中不存在的实体幻觉，或者合成违反严格医学逻辑的程序序列。纠正这些错误需要复杂且不可靠的事后过滤流程，从而损害自动生成过程的效率和可信度。

\subsection{Simulation-based Synthesis}
A third paradigm circumvents the manual annotation bottleneck entirely by leveraging graphics engines for programmatic data synthesis.
% 第三种范式完全绕过了人工标注的瓶颈，利用图形引擎进行程序化数据合成。
Pioneering works such as CLEVR \cite{CLEVR} established this paradigm, utilizing physics engines to construct structured visual reasoning benchmarks.
% CLEVR 等开创性工作确立了这种范式，它利用物理引擎构建结构化的视觉推理基准。
The distinct advantage of this approach lies in its access to the underlying simulation state—such as precise object coordinates and velocities—which guarantees absolute ground truth accuracy during dataset construction. This bypasses the noise and ambiguity inherent in extracting metadata from raw, "second-hand" video feeds.
% 这种方法的显著优势在于能够访问底层仿真状态——例如精确的物体坐标和速度——从而保证数据集构建过程中的绝对真实精度。这避免了从原始的“二手”视频流中提取元数据时固有的噪声和歧义。
Subsequently, this paradigm has been scaled to the temporal dimension, extending from static images to dynamic video. Notable frameworks like VisualRoad \cite{visualroad} leverage high-fidelity graphics engines (e.g., Unity) to render complex traffic scenarios, enabling the automated construction of large-scale synthetic video benchmarks for autonomous driving.
% 随后，这种范式扩展到了时间维度，从静态图像扩展到动态视频。诸如 VisualRoad 等知名框架利用高保真图形引擎（例如 Unity）渲染复杂的交通场景，从而能够自动构建用于自动驾驶的大规模合成视频基准。
However, while highly effective for rigid-body domains such as autonomous driving, this methodology encounters prohibitive feasibility barriers within the medical context.
% 然而，尽管这种方法在自动驾驶等刚体领域非常有效，但在医疗领域却面临着难以克服的可行性障碍。
Accurately modeling the deformable dynamics of human tissue and the fluid complexity of surgical interventions requires computational resources that scale non-linearly with realism.
% 精确模拟人体组织的可变形动力学和手术操作中的流体复杂性需要大量的计算资源，而这些资源会随着真实性的增加呈非线性增长。
Consequently, synthetic medical datasets frequently suffer from a severe "Sim-to-Real" domain shift.
% 因此，合成的医疗数据集经常面临严重的“模拟到现实”的领域偏差，导致基于这些无菌模拟训练的模型难以应对临床现实中的噪声和个体差异。
To resolve these dichotomies, our work proposes a hybrid framework that synthesizes the strengths of real-world data capture and structured logic generation.
% 为了解决这些矛盾，我们的工作提出了一种混合框架，该框架融合了真实世界数据采集和结构化逻辑生成的优势。
We leverage real-world video data to obviate the Sim-to-Real gap, while simultaneously constructing a fine-grained, cross-modal Knowledge Graph to orchestrate the query generation process.
% 我们利用真实世界的视频数据来弥合模拟到现实的差距，同时构建一个细粒度的跨模态知识图谱来协调查询生成过程。
This approach achieves the scalability of automated synthesis while strictly enforcing the visual alignment and logical determinism absent in pure LLM-based methods, effectively bridging the semantic gap without sacrificing data integrity.
% 这种方法实现了自动合成的可扩展性，同时严格执行了纯粹基于 LLM 的方法所缺乏的视觉对齐和逻辑确定性，有效地弥合了语义差距，而没有牺牲数据完整性。